% English draft based on sections/06_animism_ja.tex
\section{Animism: Installation of Stopping in the Environment}

Animism is treated here not as an error about spirits, but as an early operating form of stopping structures. It installs stopping loci in the environment.

In silent or under-responsive regions, inference cannot obtain sufficient return from the world. Animistic operation assigns stopping power to places, objects, winds, stones, rivers, animals, or events. Structurally, this means that an external element is made into a locus at which inference may stop, orient, or act.

This does not require the theory to assert the existence of spirits. Nor does it require the theory to dismiss animism as illusion. The question is structural: what operation does animism perform for finite intelligence? Its answer is that animism distributes $L_E$ into the environment.

This form is close to environmental mutual constraint. The stopping locus is not yet centralized. It is local, concrete, and embedded in direct interaction between agent and environment. In this sense, animism is an installation operation:
\[
\text{environmental locus} \;\mapsto\; E\text{-installation site}.
\]

The importance of animism in this theory is that it shows stopping before centralization. The environment itself is used as a distributed surface on which stopping can be projected.

\bigskip
